\section{Convergence theorem and proof}\label{sec:convergence}
\begin{theorem}\label{thm:convergence}
In exact real arithmetic, the sequence defined by \eqref{eq:newton} satisfies
$x_n>\rtwo$ for every integer $n\geq0$, decreases strictly, and converges to $\rtwo$.
\end{theorem}
\begin{proof}
Let $\alpha=\rtwo$. For any $x>0$, direct algebra using $\alpha^2=2$ gives
\begin{equation}\label{eq:error-identity}
\frac12\left(x+\frac{2}{x}\right)-\alpha
=\frac{x^2-2\alpha x+\alpha^2}{2x}
=\frac{(x-\alpha)^2}{2x}.
\end{equation}
Since $x_0=2>\alpha$, induction shows that each iterate is positive and strictly
above $\alpha$. In particular, the recurrence is defined at every step.
Moreover,
\[x_{n+1}-x_n=\frac{2-x_n^2}{2x_n}<0.\]
Thus the sequence is strictly decreasing and bounded below. It has a limit
$L\geq\alpha>0$. Continuity of $x\mapsto (x+2/x)/2$ on the positive real axis
allows passage to the limit in the recurrence, yielding
$L=(L+2/L)/2$. Consequently $L^2=2$, and positivity forces $L=\alpha$.
\end{proof}
The starting value is part of the statement. A positive start below the root
need not decrease on the first step; a zero start is not defined. More generally,
\eqref{eq:error-identity} puts every positive start other than the root strictly
above the root after one step. We keep $x_0=2$ fixed throughout this note to
avoid mixing this broader observation with the experiment's parameters.
